\chapter{Morphology: The Engine of Animacy}

\section{The Active-Stative Meta-System}
The PED language lacked the Nominative-Accusative framework that defines modern European and Dravidian syntax. Instead, it operated on a binary Animacy Hierarchy. The grammatical role of a noun was determined not by its relationship to the verb, but by its inherent capacity for volition (will).

\subsection{The Active Class (*-S)}
Agential beings capable of independent motion—humans, large animals, fire, and wind—were marked with the \textit{*-S} suffix.
\begin{itemize}
    \item \textbf{Vedic Reflex:} Survived as the Masculine/Feminine nominative singular \textit{*-s}.
    \item \textbf{Old Tamil Reflex:} Survived as the 'Rational' (uyartiṇai) markers.
\end{itemize}

\subsection{The Inactive Class (*-N)}
Passive objects or experiencers of a state were marked with the \textit{*-N} suffix.
\begin{itemize}
    \item \textbf{Vedic Reflex:} Survived as the Neuter nominative/accusative \textit{*-m}.
    \item \textbf{Old Tamil Reflex:} Survived as the 'Irrational' (akṟiṇai) neuter singular \textit{*-n/*-m}.
\end{itemize}

\section{The Noun Declension Paradigm}
The PED noun was a modular unit capable of carrying five primary case/state markers. The following table maps these suffixes to their most ancient attested daughter reflexes.

\begin{longtable}{@{}llll@{}}
\toprule
\textbf{Case/State} & \textbf{PED Suffix} & \textbf{Vedic Reflex} & \textbf{Tamil Reflex} \\ \midrule
\endfirsthead
\toprule
\textbf{Case/State} & \textbf{PED Suffix} & \textbf{Vedic Reflex} & \textbf{Tamil Reflex} \\ \midrule
\endhead
\bottomrule
\endfoot
\bottomrule
\lastfoot
Active/Agent & *-S & Nom. *-s & Rational Marker \\
Inactive/Target & *-N & Neut. *-m & Accusative *-n \\
Directional & *-K & (and) *-kʷe & Dative *-ku \\
Locative & *-R & Loc. *-er & Locative *-il \\
Instrumental & *-P & Plural *-bhi & Increment *-p- \\
Reflexive & *-v- & (Middle Voice) & koḷ- (Take) \\
\end{longtable}

\section{The Verbal Complex and the Agreement Engine}
PED agreement was strictly deictic, mapping the relationship between the self and the other.

\subsection{Primitive Agreement Markers}
\begin{table}[h]
\centering
\begin{tabular}{@{}llll@{}}
\toprule
\textbf{Person} & \textbf{PED Primitive} & \textbf{Vedic Development} & \textbf{Tamil Development} \\ \midrule
1st (Self) & *eK (Active) & *e-ǵo $\rightarrow$ *-mi & *-ē\d{n} \\
2nd (Other) & *tV (Deictic) & *tu- $\rightarrow$ *-si & *-āy \\
3rd (That) & *sV (Anaphor) & *so- $\rightarrow$ *-ti & *-tu \\ \bottomrule
\end{tabular}
\caption{The Agreement Matrix.}
\end{table}

\section{The 'r/n' Heteroclitic Fossil}
One of the most profound 'glitches' in the PED mother tongue was the heteroclitic declension, where archaic nouns for natural forces alternated between an \textit{*r}-stem (Active/Dynamic) and an \textit{*n}-stem (Inactive/Resultative).
\begin{center}
\textit{*wódr̥} (Running Water) vs. \textit{*wédn̥s} (The Result of Drinking).
\end{center}
We identify this same alternation in the Old Tamil first-person pronouns: \textit{yān} (nominative/static) vs. \textit{en-} (oblique/targeted).

\section{The 'Interiority' Nasal cluster (*M-N)}
A dedicated morphological cluster governed the 'Interior Life'. The nasal combination \textit{*M-N} was restricted to verbs of thinking, staying, and the concept of the 'self'.
\begin{table}[h]
\centering
\begin{tabular}{lll} \toprule
\textbf{Language} & \textbf{Reflex} & \textbf{Meaning} \\ \midrule
Vedic & \textit{manas} & Mind/Spirit \\
Old Tamil & \textit{maṉam} & Heart/Will \\
Old Tibetan & \textit{man} & To stay/To dwell \\
PED Root & \textit{*M-e-N-} & The Staying Mind \\ \bottomrule
\end{tabular}
\caption{The Interiority Cluster.}
\end{table}

\section{Transitive Gemination (Consonant Hardening)}
PED indicated the transition from an agential state to a causal impact by hardening the root consonant. This is the origin of the Dravidian gemination rules (e.g., \textit{*āṭu} "to dance" vs. \textit{*āṭṭu} "to make dance"). In PIE, this was often handled by the \textit{s-mobile} prefix, which added agential 'friction' to the root.
